\documentclass[11pt]{article}
\usepackage{amsmath,amssymb,amsthm}
\usepackage[margin=1.1in]{geometry}
\usepackage{booktabs}
\usepackage{hyperref}

\newtheorem{remark}{Remark}

\title{Scaling at Desk Size:\\
Which Resource Is the Bottleneck, Measured}
\author{Yuval Lavie}
\date{\today}

\begin{document}
\maketitle

\begin{abstract}
Rung 4 chose 4 layers and $d = 256$ by taste. This note replaces
taste with two sweeps under a fixed budget (3000 steps of the rung-4
recipe, same seed, same validation slice): model size from $0.06$M to
$7.5$M parameters at full data, and data from $10\%$ to $100\%$ of
the corpus at fixed model size. The two curves answer the resource
question unambiguously: parameter gains shrink to nothing while the
train--val gap widens, and every data doubling keeps paying --- at
$10^6$ characters, the next nat lives in more text, not more model.
The Chinchilla lesson, reproduced on one desk in fifteen minutes.
\texttt{scaling\_experiments.py} (results also in
\texttt{scaling\_results.json}).
\end{abstract}

\section{The two sweeps}

\begin{center}
\begin{tabular}{lrrrr}
\toprule
\multicolumn{5}{c}{\textbf{Param sweep} (full data, 3000 steps)}\\
model & params & train NLL & val NLL & gap\\
\midrule
1 layer, $d{=}64$ & $0.06$M & $2.081$ & $2.147$ & $0.066$\\
2 layers, $d{=}128$ & $0.42$M & $1.616$ & $1.780$ & $0.164$\\
3 layers, $d{=}192$ & $1.37$M & $1.360$ & $1.557$ & $0.197$\\
4 layers, $d{=}256$ & $3.21$M & $1.227$ & $1.475$ & $0.248$\\
6 layers, $d{=}320$ & $7.46$M & $1.126$ & $1.459$ & $0.332$\\
\bottomrule
\end{tabular}
\end{center}

\begin{center}
\begin{tabular}{lrrr}
\toprule
\multicolumn{4}{c}{\textbf{Data sweep} (4L/$d$256, 3000 steps)}\\
train data & train NLL & val NLL & \\
\midrule
$10\%$ ($\sim$100k chars) & $0.134$ & $3.823$ &
  \emph{memorized; worse than a bigram}\\
$25\%$ & $0.778$ & $2.272$ & \\
$50\%$ & $1.109$ & $1.863$ & \\
$100\%$ & $1.227$ & $\mathbf{1.475}$ & \\
\bottomrule
\end{tabular}
\end{center}

\section{Reading the curves}

\paragraph{Params stall.} Each model-size step multiplies parameters
by $\sim$2--7$\times$; the validation gains are $0.37$, $0.22$,
$0.08$, and finally $0.017$ nats --- the last $4.3$M parameters
bought almost nothing. Meanwhile the train--val gap grows
monotonically, $0.066 \to 0.332$: the big models have plenty of
capacity and use it to fit noise. They are not starved of parameters;
they are starved of data.

\paragraph{Data keeps paying.} With the model fixed, every data
increase improves validation, and the steps are enormous
($3.82 \to 2.27 \to 1.86 \to 1.48$). The $10\%$ row is the
experiment's best teacher: $3.2$M parameters against $10^5$
characters reaches train NLL $0.134$ --- essentially a lossless
memory of the training text --- while validation lands at $3.82$,
\emph{worse than rung 1's bigram} ($2.48$). A model that can quote
its training set perfectly and cannot predict fresh text is the
overfitting story of Theory/learning-theory, staged at scale: the
count-table pathology of rung 2's U-turn, now happening to a
transformer.

\paragraph{The verdict.}
\begin{equation}
  \boxed{\;\text{At } 10^6 \text{ chars: params } 2.147 \to 1.459
  \text{ (stalling)}, \quad \text{data } 3.823 \to 1.475
  \text{ (still paying)}\;}
  \label{eq:verdict}
\end{equation}
The next improvement to the ladder's GPT is a bigger corpus, not a
bigger model --- which is exactly the Chinchilla result (compute is
best spent balancing model \emph{and} data, and most models of this
era were data-starved), measured with nine desk-scale runs. All three
shape claims --- diminishing param gains, growing overfit gap,
monotone data gains --- are asserts in the script, not prose.

\begin{remark}[What the fixed budget hides]
Both sweeps hold the step budget constant, so ``more parameters''
also means ``more compute per step'' but \emph{not} more optimization
steps; a full scaling-law study would trade these against each other
(loss vs compute at optimal allocation). At this corpus size the
conclusion would not change: the data curve dominates every other
axis. The honest scope of \eqref{eq:verdict} is: at fixed budget, on
this corpus.
\end{remark}

\begin{remark}[Next]
The measured bottleneck names the next experiment: a larger corpus
(more Shakespeare-adjacent drama, or the full Gutenberg Shakespeare)
with the same ladder discipline --- and, on the model side, the speed
knobs (bf16 autocast, flash attention) that buy steps rather than
parameters.
\end{remark}

\end{document}
